\section{Background}
\label{sec:background}

%Note that this section does not contain any information that refers to the number of studies etc. 

% Name subsection something like this: The problem of customer churn and AI powered marketing retention strategies
%\DD{}{Add a preamble here, summarizing what this Section is intended/serves for}

This section is divided into two parts. The first subsection outlines the fundamental concepts of customer churn, distinguishes between its various types, and traces the progression of churn management strategies. It highlights how the increasing availability of customer data and advancements in predictive analytics have transformed retention practices from more generalized retention marketing activities to more targeted approaches based on predicted risk of customer churn.

The second subsection presents an overview of the machine learning methods employed in the selected studies analyzed in this SLR, providing a reference point to help readers better contextualize and understand the subsequent sections of the document.
 
\subsection{\textcolor{blue}{The customer churn problem and mitigation approaches}}

\textcolor{blue}{As stated by \cite{S37-ID826-Buckinx}, ``Churn analysis typically tries to define predictors of customer defection.'' Customer churn refers to the situation in which a customer terminates or substantially reduces their relationship with a firm, for example, by closing an account, switching to another provider, or no longer purchasing its goods or services. This behavior can occur in both contractual and non-contractual settings. In contractual industries, such as telecommunications or insurance, churn typically occurs when a customer cancels or fails to renew a formal service agreement, and is therefore observed as a form of total defection. In non-contractual contexts, customers can continuously adjust their purchase behavior without explicitly terminating the relationship. In these settings, churn is often studied through changes in purchasing patterns and partial defection, defines as a sustained reduction in purchase frequency or spending with the focal retailer that may eventually lead to complete defection.}


Churn is commonly categorized into two main types: involuntary and voluntary \cite{park2022not}. Involuntary churn arises from factors outside the customer’s control, such as relocation to an area not served by the business, financial hardship, or major life events like illness or death. Voluntary churn, on the other hand, occurs when a customer consciously decides to disengage. This may be driven by dissatisfaction with product or service quality, negative customer experiences, more competitive pricing or offerings from other providers, or a general perception of declining value.

While involuntary churn is often difficult for businesses to influence, minimizing voluntary churn is of strategic importance, particularly given that acquiring new customers is significantly more expensive than retaining existing ones \cite{ahmad2001customer}. Consequently, many businesses prioritize customer retention strategies and increasingly leverage churn prediction models to identify high-risk customers. One widespread approach was the introduction of loyalty programs, designed to raise the barrier to switching and encourage engagement \cite{lewis2004influence}. These programs also gave retailers access to valuable data about customer behavior, enabling deeper insights into purchasing patterns \cite{stourm2020refocusing}.

With access to detailed customer data, businesses adopted more sophisticated strategies by leveraging historical customer behavior to segment individuals based on their predicted churn risk. This segmentation enabled the shift from broad marketing campaigns to more targeted efforts. Marketing budgets were increasingly directed toward customers identified as having a high risk of churn \cite{lemmens2020managing}, allowing for the deployment of tailored interventions, such as personalized offers, loyalty incentives, or service enhancements to foster continued engagement. In this context, accurate and early churn prediction emerged as a critical tool for optimizing resource allocation and improving customer retention outcomes.

Machine learning (ML) techniques have been applied to the task of churn prediction over the years \cite{churn-prediction-survery-manzoor-2024}. ML involves the use of algorithms that learn patterns from data without explicit programming, allowing models to capture complex, often nonlinear relationships. The inherent complexity of churn prediction, coupled with the large volumes of data typically available in business contexts, makes this task particularly well-suited to ML approaches. These techniques have been employed prevalently in contractual sectors including banking and telecommunications but less frequently in non-contractual environments such as retail.

In the following section, we provide a structured overview of the AI and ML techniques employed across the selected studies, with a particular emphasis on their use in churn prediction within the retail context. 

\subsection{\textcolor{blue}{Machine learning techniques for churn prediction in the retail sector}}

\textcolor{blue}{Considering the large datasets available to retailers, ML is particularly well suited for churn prediction. Although our SLR considered primary studies employing any AI technique (e.g., deep learning, reinforcement learning), all studies that met our relevance and quality criteria ultimately applied ML methods.} ML, a subfield of AI, focuses on developing models capable of making predictions or decisions based on data-driven learning \cite{Sarker2021}. It encompasses a broad spectrum of techniques, each tailored to different problem domains and data characteristics. This section presents an overview and definitions of the machine learning methods identified in the selected studies included in this SLR, utilized for the task of customer churn prediction.

%This serves to provide a common understaning of the machine learning techniques to better allow the reader to comprehend the results and analysis performed in the paper. 
\subsubsection{Tree-Based Methods}

Tree-based methods utilize decision trees as their foundational structure, offering interpretability and the ability to handle both numerical and categorical data \cite{james2023tree}. \textit{Decision Trees} are non-parametric supervised learning algorithms used for classification and regression tasks. They work by recursively splitting the dataset into subsets based on feature values, forming a tree-like structure of decisions. \textit{Random Forests} are ensemble methods that construct multiple decision trees during training and output the mode (for classification) or mean (for regression) of the predictions made by individual trees. \textit{Extra Trees}, or Extremely Randomized Trees, introduce additional randomness by choosing split points randomly, which can lead to better generalization and reduced overfitting. \textit{Boosting classifiers}, including \textit{XGBoost}, \textit{CatBoost}, \textit{AdaBoost}, and \textit{LogitBoost}, build strong learners by combining multiple weak learners, typically decision trees. \textit{XGBoost} is known for its efficiency and performance, particularly on structured data, while \textit{CatBoost} handles categorical variables natively, reducing preprocessing efforts. \textit{AdaBoost} improves classification by adjusting weights to emphasize previously misclassified instances. \textit{LogitBoost} optimizes a logistic loss function through additive modeling. \textit{Bagging classifiers}, such as Bootstrap Aggregating, train multiple models on different bootstrap samples of the data and combine their outputs to improve stability and accuracy. \textcolor{blue}{Tree based methods demonstrate strong performance when trained on an imbalanced dataset. For the churn prediction problem, where the ratio of churners is significantly lower than that of non-churners, this characteristic of tree-based methods would be advantageous \cite{breiman2001random,chen2016xgboost}. }
\subsubsection{Distance-Based Methods}

Distance-based methods rely on geometric interpretations in feature space \cite{hsu2002comparison}. \textit{Support Vector Machines (SVMs)} construct optimal hyperplanes to separate classes with maximum margin and are effective in high-dimensional spaces. \textit{K-Nearest Neighbors (KNN)} classifies data points based on the majority class among the 'k' closest examples in the training set, offering simplicity and interpretability. \textit{Multivariate Adaptive Regression Splines (MARS)} \cite{friedman1995introduction} is a flexible regression technique that captures nonlinear relationships by fitting piecewise linear regressions, making it suitable for modeling complex interactions between variables \textcolor{blue}{such as with churn prediction.}

\subsubsection{Linear and Discriminant Analysis}

The methods in this category are often classification and regression tasks. \textit{Ridge Regression} is a technique that adds L2 regularization to linear regression, penalizing large coefficients to prevent overfitting. Linear Discriminant Analysis and Quadratic Discriminant Analysis are both derive a discriminant function for each class and classify observations based on the function's output \cite{tharwat2016linear}. \textit{Linear Discriminant Analysis (LDA)} is a method that finds a linear combination of features that best separates two or more classes. \textit{Quadratic Discriminant Analysis (QDA)} is similar to LDA but allows for each class to have its own covariance matrix, capturing more complex class distributions. \textit{Logistic Regression} is a statistical model that predicts the probability of a binary outcome using a logistic function \cite{hosmer2013applied}. \textit{Boosted Logistic Regression} enhances logistic regression by applying boosting techniques to improve predictive performance. \textcolor{blue}{These methods are perhaps simplistic for the churn prediction problem at hand, due to their limited ability to generalize relationships.}

\subsubsection{Stochastic} 
Stochastic methods are probabilistic models that make predictions based on the likelihood of outcomes, incorporating randomness in their processes. \textit{Naive Bayes} \cite{gladence2015statistical} includes classification techniques based on Bayes' Theorem, assuming independence between predictors. \textit{Poisson Processes} \cite{kingman1992poisson} are stochastic processes that model the occurrence of events over time, often used for count data. \textcolor{blue}{Although such methods can be effective for the right task, the assumptions made in these probabilistic models might not apply to the churn prediction problem, where there can be complex and dependent interrelations between the feature variables \cite{sahami1996learning}.} 

\subsubsection{Neural Networks}

Neural networks are inspired by the structure and functioning of the human brain and consist of layers of interconnected nodes (neurons) that process input data \cite{janiesch2021machine}. \textit{Artificial Neural Networks (ANNs)} are composed of multiple layers of neurons and are used for a wide range of prediction tasks. \textit{Recurrent Neural Networks (RNNs)} are designed to process sequential data by incorporating recurrent connections, where the output from one time step is fed as input to the next. This structure enables RNNs to capture temporal dependencies in data such as time series or text. \textit{Long Short-Term Memory (LSTM)} networks, a specialized form of RNN, address the limitations of standard RNNs by incorporating memory cells and gating mechanisms, allowing them to learn long-term dependencies more effectively. \textit{Convolutional Neural Networks (CNNs)}, typically used in image processing tasks, consist of convolutional layers that automatically learn spatial hierarchies of features. \textit{Restricted Boltzmann Machines (RBMs)} are generative stochastic neural networks that learn probability distributions over input data. They are particularly useful for dimensionality reduction, feature learning, and collaborative filtering. \textcolor{blue}{Neural networks can have great generalization power due to their ability to model complex relationships, however, this generalizing ability comes at the cost of requiring a large corpus of data to extract the best performance form such models. For the churn prediction problem, retailers naturally have access to data at this scale; therefore, neural network-based approaches could be well suited \cite{ajiboye2015evaluating}. }


